\subsection{Multi-bit robustness depends on the perturbation starting measure}
The single-bit fragility index averages perturbations over every state in a basin. This measure is mathematically well-defined but differs from asking how stable an established fate is. We tested that distinction by enumerating the complete 12-node floral network state space and every Hamming-radius perturbation. We retain the published synchronous threshold update, rather than treating this experiment as a reconstruction of molecular fluctuations.

Let $a(x)$ denote the eventual attractor of state $x$, $x_f$ a fixed point, and $B_f=\{x:a(x)=a(x_f)\}$ its basin. Two retention functions are
\[
R_f^{\mathrm{fixed}}(k)=\frac{1}{\binom{12}{k}}\sum_{\delta:|\delta|=k}\mathbf{1}[a(x_f\mathbin{\mathrm{xor}}\delta)=a(x_f)]
\]
and
\[
R_f^{\mathrm{basin}}(k)=\frac{1}{|B_f|\binom{12}{k}}\sum_{x\in B_f}\sum_{\delta:|\delta|=k}\mathbf{1}[a(x\mathbin{\mathrm{xor}}\delta)=a(x_f)].
\]
The earlier fragility index equals $1-R_f^{\mathrm{basin}}(1)$. It does not equal $1-R_f^{\mathrm{fixed}}(1)$ in general. The first function models perturbing an established attractor; the second models a uniformly sampled trajectory-starting state that will eventually reach it. Neither uniform measure has been calibrated to the distribution of actual developmental states.

\begin{center}\begin{tabular}{lrrrr}\hline
Fate & Fixed $R(1)$ & Basin $R(1)$ & Fixed $q_{1/2}$ & Basin $q_{1/2}$\\\hline
sepal & 0.500 & 0.534 & 0.116 & 0.122 \\
petal & 0.417 & 0.601 & 0.105 & 0.141 \\
carpel & 0.333 & 0.319 & 0.084 & 0.083 \\
stamen & 0.167 & 0.208 & 0.068 & 0.071 \\
inflorescence & 0.667 & 0.653 & 0.167 & 0.163 \\
mutant & 0.583 & 0.681 & 0.146 & 0.173 \\
\hline\end{tabular}\end{center}
Here $q_{1/2}$ is the first tested independent per-node flip probability whose same-attractor retention falls to one half. It is obtained by weighting the exact-radius curves with the binomial distribution,
\[
P_f(q)=\sum_{k=0}^{12}\binom{12}{k}q^k(1-q)^{12-k}R_f(k).
\]
The grid resolution is $0.001$, so the quoted crossing is a grid bracket, not a fitted biological threshold. At $q=0$, retention is one. At $q=0.5$, every perturbed state is uniformly distributed across the full cube, so retention must equal $|B_f|/4096$ independently of the starting state. This is an exact consistency check on both measures. Every one-bit basin result also recovers the previously archived fragility value within its rounding precision.

The petal-versus-sepal ordering reverses. From their established fixed points, sepal retains fate under one bit flip with probability $0.500$, while petal retains it with probability $0.417$. Averaging uniformly over each basin instead gives $0.534$ for sepal and $0.601$ for petal. Petal is therefore less stable under the fixed-point measure but more stable under the basin-start measure. One scalar fragility ranking cannot represent both questions.

This is a positive methodological limit: robustness is a property of the dynamics together with a perturbation kernel and a starting-state measure. It is not a topology-only invariant, and a cross-organism comparison must specify those choices. The observed ordering reversal is exact within the transcribed model and does not rely on a noisy simulation estimate. It remains model-bound: artificial bit noise, synchronous update and uniform basin weights are not direct molecular interventions. A useful next test would weight states by measured developmental trajectories and compare update schedules under the same kernel, rather than calling the current values intrinsic canalization.

The full radius curves, independent-flip grid values and source-code checksum are in \path{results/flower_perturbation_radius.json}; the executable audit is \path{experiments/flower_perturbation_radius.py}. No new empirical biological discovery or causal topology law is claimed.
